% GENERATED FILE. Edit canonical lessons, then run npm run generate:books.
% canonical-source-hash: fnv1a64:e4b60541f5c2386c
% canonical-lessons: JA-C414-shitsudo, JA-C414-jimejime, JA-C414-hizashi, JA-C414-hikage, JA-C414-mabushii

\chapter{Hot, Damp Days}
\label{ch:ja-hot-damp-days}
\hlchaptermodality{414}

\begin{chapteropening}
\textbf{By the end of this chapter:} \emph{I can talk about heat, damp and sunshine on a trip.}

Complete the last lesson of chapter 414: I can talk about heat, damp and sunshine on a trip.
\end{chapteropening}

\section[\texorpdfstring{\ja{しつど} (shitsudo)}{しつど (shitsudo)}]{\ja{しつど} (shitsudo) — humidity}

\label{lesson:JA-C414-shitsudo}



\begin{quote}
Before the new one: say the Japanese for a continent, then the Japanese for a desert.
\end{quote}

\subsection*{You'll want to know: \ja{しつど}}
\textbf{\ja{しつど}} — \emph{shitsudo} — \textquotedblleft{}humidity\textquotedblright{}.

\begin{quote}
\textbf{\ja{なつの} \ja{日本}\ja{は} \ja{しつどが} \ja{たかいです。}} — \emph{natsu no nihon wa shitsudo ga takai desu} — \textquotedblleft{}Japan is humid in summer.\textquotedblright{}
\end{quote}

\begin{cousinweb}
\textbf{\ja{しつど}} (\emph{shitsudo}) is Sino-Japanese (kango): two Chinese characters read with their on'yomi, readings borrowed from Chinese. They are \emph{shitsu}, \textquotedblleft{}damp\textquotedblright{} (Middle Chinese \emph{syip}) and \emph{do}, \textquotedblleft{}degrees\textquotedblright{} (\emph{duH}) (KANJIDIC2; Middle Chinese in Baxter's transcription).
\end{cousinweb}

\begin{grammarlens}[title={What you've built}]
One more word for getting around.
\end{grammarlens}

\subsection*{Your turn}
\begin{itemize}
\raggedright
  \item \emph{Say it:} \emph{shitsudo}
  \item \emph{Say it:} \emph{shitsudo}, once more
  \item \emph{Say it:} \emph{natsu no nihon wa shitsudo ga takai desu} — Japan is humid in summer.
  \item \emph{Recall:} say the Japanese for a continent, then the Japanese for a desert, then say \emph{shitsudo} again
\end{itemize}

\begin{tcolorbox}[breakable,colback=teal!4,colframe=teal!35!black,title={Before you move on}]
What does \ja{しつど} mean? (\textquotedblleft{}Humidity\textquotedblright{}.) Say it once more.
\end{tcolorbox}
\section[\texorpdfstring{\ja{じめじめ} (jimejime)}{じめじめ (jimejime)}]{\ja{じめじめ} (jimejime) — damp and humid}

\label{lesson:JA-C414-jimejime}



\begin{quote}
Before the new one: say the Japanese for a desert, then the Japanese for humidity.
\end{quote}

\subsection*{You'll want to know: \ja{じめじめ}}
\textbf{\ja{じめじめ}} — \emph{jimejime} — \textquotedblleft{}damp and humid\textquotedblright{}.

\begin{quote}
\textbf{\ja{つゆは} \ja{まいにち} \ja{じめじめしています。}} — \emph{tsuyu wa mainichi jimejime shite imasu} — \textquotedblleft{}In the rainy season it is damp and sticky every day.\textquotedblright{}
\end{quote}

\begin{cousinweb}
\textbf{\ja{じめじめ}} (\emph{jimejime}) is a native mimetic word that pictures damp, sticky air: an imitative Japanese form, inherited rather than borrowed, whose older origin is not settled (JMdict).
\end{cousinweb}

\begin{grammarlens}[title={What you've built}]
One more word for getting around.
\end{grammarlens}

\subsection*{Your turn}
\begin{itemize}
\raggedright
  \item \emph{Say it:} \emph{jimejime}
  \item \emph{Say it:} \emph{jimejime}, once more
  \item \emph{Say it:} \emph{tsuyu wa mainichi jimejime shite imasu} — In the rainy season it is damp and sticky every day.
  \item \emph{Recall:} say the Japanese for a desert, then the Japanese for humidity, then say \emph{jimejime} again
\end{itemize}

\begin{tcolorbox}[breakable,colback=teal!4,colframe=teal!35!black,title={Before you move on}]
What does \ja{じめじめ} mean? (\textquotedblleft{}Damp and humid\textquotedblright{}.) Say it once more.
\end{tcolorbox}
\section[\texorpdfstring{\ja{ひざし} (hizashi)}{ひざし (hizashi)}]{\ja{ひざし} (hizashi) — sunlight; the sun's rays}

\label{lesson:JA-C414-hizashi}



\begin{quote}
Before the new one: say the Japanese for humidity, then the Japanese for damp and humid.
\end{quote}

\subsection*{You'll want to know: \ja{ひざし}}
\textbf{\ja{ひざし}} — \emph{hizashi} — \textquotedblleft{}sunlight; the sun's rays\textquotedblright{}.

\begin{quote}
\textbf{\ja{なつの} \ja{ひざしは} \ja{つよいです。}} — \emph{natsu no hizashi wa tsuyoi desu} — \textquotedblleft{}The summer sun is strong.\textquotedblright{}
\end{quote}

\begin{cousinweb}
\textbf{\ja{ひざし}} (\emph{hizashi}) is a native Japanese word (wago), inherited rather than borrowed from Chinese, built from \emph{hi}, \textquotedblleft{}sun\textquotedblright{} and \emph{zashi}, \textquotedblleft{}shining\textquotedblright{} (from \emph{sasu}); its kanji were attached later as kun'yomi (JMdict).
\end{cousinweb}

\begin{grammarlens}[title={What you've built}]
One more word for getting around.
\end{grammarlens}

\subsection*{Your turn}
\begin{itemize}
\raggedright
  \item \emph{Say it:} \emph{hizashi}
  \item \emph{Say it:} \emph{hizashi}, once more
  \item \emph{Say it:} \emph{natsu no hizashi wa tsuyoi desu} — The summer sun is strong.
  \item \emph{Recall:} say the Japanese for humidity, then the Japanese for damp and humid, then say \emph{hizashi} again
\end{itemize}

\begin{tcolorbox}[breakable,colback=teal!4,colframe=teal!35!black,title={Before you move on}]
What does \ja{ひざし} mean? (\textquotedblleft{}Sunlight; the sun's rays\textquotedblright{}.) Say it once more.
\end{tcolorbox}
\section[\texorpdfstring{\ja{ひかげ} (hikage)}{ひかげ (hikage)}]{\ja{ひかげ} (hikage) — shade (out of the sun)}

\label{lesson:JA-C414-hikage}



\begin{quote}
Before the new one: say the Japanese for damp and humid, then the Japanese for the sun's rays.
\end{quote}

\subsection*{You'll want to know: \ja{ひかげ}}
\textbf{\ja{ひかげ}} — \emph{hikage} — \textquotedblleft{}shade (out of the sun)\textquotedblright{}.

\begin{quote}
\textbf{\ja{ひかげで} \ja{すこし} \ja{やすみました。}} — \emph{hikage de sukoshi yasumimashita} — \textquotedblleft{}I rested a little in the shade.\textquotedblright{}
\end{quote}

\begin{cousinweb}
\textbf{\ja{ひかげ}} (\emph{hikage}) is a native Japanese word (wago), inherited rather than borrowed from Chinese, built from \emph{hi}, \textquotedblleft{}sun\textquotedblright{} and \emph{kage}, \textquotedblleft{}shadow\textquotedblright{}; its kanji were attached later as kun'yomi (JMdict).
\end{cousinweb}

\begin{grammarlens}[title={What you've built}]
One more word for getting around.
\end{grammarlens}

\subsection*{Your turn}
\begin{itemize}
\raggedright
  \item \emph{Say it:} \emph{hikage}
  \item \emph{Say it:} \emph{hikage}, once more
  \item \emph{Say it:} \emph{hikage de sukoshi yasumimashita} — I rested a little in the shade.
  \item \emph{Recall:} say the Japanese for damp and humid, then the Japanese for the sun's rays, then say \emph{hikage} again
\end{itemize}

\begin{tcolorbox}[breakable,colback=teal!4,colframe=teal!35!black,title={Before you move on}]
What does \ja{ひかげ} mean? (\textquotedblleft{}Shade (out of the sun)\textquotedblright{}.) Say it once more.
\end{tcolorbox}
\section[\texorpdfstring{\ja{まぶしい} (mabushii)}{まぶしい (mabushii)}]{\ja{まぶしい} (mabushii) — dazzling; too bright}

\label{lesson:JA-C414-mabushii}



\begin{quote}
Before the new one: say the Japanese for the sun's rays, then the Japanese for shade.
\end{quote}

\subsection*{You'll want to know: \ja{まぶしい}}
\textbf{\ja{まぶしい}} — \emph{mabushii} — \textquotedblleft{}dazzling; too bright\textquotedblright{}.

\begin{quote}
\textbf{\ja{ゆきが} \ja{ひかって、まぶしいです。}} — \emph{yuki ga hikatte, mabushii desu} — \textquotedblleft{}The snow is shining and it is dazzling.\textquotedblright{}
\end{quote}

An i-adjective.

\begin{cousinweb}
\textbf{\ja{まぶしい}} (\emph{mabushii}) is a native Japanese word (wago), inherited rather than borrowed from Chinese; its kanji was attached later as a kun'yomi (JMdict).
\end{cousinweb}

\begin{grammarlens}[title={What you've built}]
One more word for getting around.
\end{grammarlens}

\subsection*{Your turn}
\begin{itemize}
\raggedright
  \item \emph{Say it:} \emph{mabushii}
  \item \emph{Say it:} \emph{mabushii}, once more
  \item \emph{Say it:} \emph{yuki ga hikatte, mabushii desu} — The snow is shining and it is dazzling.
  \item \emph{Recall:} say the Japanese for the sun's rays, then the Japanese for shade, then say \emph{mabushii} again
\end{itemize}

\begin{tcolorbox}[breakable,colback=teal!4,colframe=teal!35!black,title={Before you move on}]
What does \ja{まぶしい} mean? (\textquotedblleft{}Dazzling; too bright\textquotedblright{}.) Say it once more.
\end{tcolorbox}
